% year: תשפ"ו
% semester: א
% moed: א
% number: 3
% points: 16
% categories: taylor
% official_solution: yes
% source: exams/מבחנים/תשפו/AnyScanner_02_01_2026(1) (2).pdf

\begin{question}
(16 נק') חשבו בעזרת קירוב טיילור מסדר שני את $\ln\left(\frac{6}{5}\right) = \ln(1.2)$, והעריכו את השגיאה.

שימו לב שהקירוב וחסם השגיאה צריכים להיות מספרים רציונליים.
\end{question}

\begin{hint}
פתחו את $f(x) = \ln x$ סביב $x_0 = 1$ (שם הערכים של $f$ ונגזרותיה רציונליים) והשתמשו בשארית לגרנז' מסדר שני.
\end{hint}

\begin{solution}
נגדיר $f(x) = \ln x$ ונפתח סביב $x_0 = 1$, נקודה קרובה ל-$1.2$ שבה כל הערכים ידועים. $f$ גזירה אינסוף פעמים ב-$(0,\infty)$:
\[ f(x) = \ln x,\quad f'(x) = \frac{1}{x},\quad f''(x) = -\frac{1}{x^2},\quad f'''(x) = \frac{2}{x^3}, \]
ולכן $f(1) = 0$, $f'(1) = 1$, $f''(1) = -1$.

\textbf{פולינום טיילור מסדר שני סביב $1$:}
\[ P_2(x) = f(1) + f'(1)(x - 1) + \frac{f''(1)}{2!}(x - 1)^2 = (x - 1) - \frac{(x - 1)^2}{2}. \]
לכן
\[ \ln(1.2) \approx P_2(1.2) = 0.2 - \frac{0.04}{2} = 0.18 = \frac{9}{50}. \]

\textbf{הערכת השגיאה.} לפי משפט טיילור עם שארית לגרנז', קיימת נקודה $c$ בין $1$ ל-$1.2$ כך ש-
\[ R_2(1.2) = \ln(1.2) - P_2(1.2) = \frac{f'''(c)}{3!}(1.2 - 1)^3 = \frac{2}{6c^3}\cdot (0.2)^3 = \frac{0.008}{3c^3}. \]
מכיוון ש-$1 < c < 1.2$, מתקיים $c^3 > 1$, ולכן $\frac{1}{c^3} < 1$ ו-
\[ 0 < R_2(1.2) < \frac{0.008}{3} = \frac{8}{3000} = \frac{1}{375}. \]

\textbf{תשובה:} $\ln(1.2) \approx 0.18 = \frac{9}{50}$, והשגיאה קטנה מ-$\frac{1}{375}$ (בפרט, מכיוון שהשארית חיובית, $0.18 < \ln(1.2) < 0.18 + \frac{1}{375}$).

(בדיקה: $\ln 1.2 \approx 0.18232$, והשגיאה בפועל היא כ-$0.0023 < \frac{1}{375} \approx 0.00267$.)
\end{solution}
